% Generated draft: requires mathematical and source review.
\documentclass[11pt]{article}
\usepackage[margin=25mm]{geometry}
\usepackage{fontspec}
\setmainfont{Latin Modern Roman}
\usepackage{amsmath,amssymb,mathtools,unicode-math}
\setmathfont{Latin Modern Math}
\usepackage{enumitem,needspace,xurl,hyperref}
\setlength{\parindent}{0pt}
\setlength{\parskip}{5pt}
\setlength{\emergencystretch}{3em}
\newcommand{\N}{\mathbb{N}}
\newcommand{\Z}{\mathbb{Z}}
\newcommand{\Q}{\mathbb{Q}}
\newcommand{\R}{\mathbb{R}}
\newcommand{\C}{\mathbb{C}}
\newcommand{\F}{\mathbb{F}}
\newcommand{\A}{\mathbb{A}}
\newcommand{\PP}{\mathbb{P}}
\newcommand{\E}{\mathbb{E}}
\newcommand{\eps}{\varepsilon}
\newcommand{\dd}{\,\mathrm{d}}
\begin{document}
\section*{UnsolvedMath ID: 800003}
\textbf{Source problem code:} catalogue entry
\subsection{Definitions and mathematical statement}
A semigroup is a nonempty set equipped with an associative binary operation. A group \(G\) is periodic if every \(g\in G\) has finite order, that is, \(\forall g\in G\,\exists n\geq 1\colon g^n=1_G\). For subsets \(A,B\subseteq G\), define their setwise product by
\[
AB=\{ab:a\in A,\ b\in B\}.
\]
Two conventions for the power semigroup are
\[
\mathcal P(G)=\{A:A\subseteq G\},\qquad
\mathcal P^{+}(G)=\mathcal P(G)\setminus\{\varnothing\},
\]
with setwise multiplication. An embedding of a semigroup \(S\) into either of these semigroups is an injective map \(\iota\) satisfying
\[
\iota(st)=\iota(s)\iota(t)\qquad(s,t\in S).
\]
The supplied source does not specify whether the empty subset is admitted. For an infinite periodic group, it also does not specify whether arbitrary subsets or only finite subsets are allowed. Consequently, the formulations below are conditional on the full-subset interpretation, with the empty-subset convention still to be selected from the source's Chapter 5. If only finite subsets are intended, the corresponding family of finite subsets must replace the target power semigroup.

\textbf{Problem (S. Bershadsky and S. Kublanovsky).} Fix the intended power-semigroup convention, denoting its target by \(\mathcal T(G)\). Given an arbitrary finite semigroup \(S\), specified by its multiplication table, does there exist an algorithm that always terminates and correctly decides whether
\[
\exists\text{ a finite group }G\quad\exists\text{ an injective semigroup homomorphism }S\longrightarrow\mathcal T(G)?
\]
The parenthetical wording ``periodic group'' gives a second apparent decision question: does there exist an algorithm that always terminates and correctly decides, for every such input \(S\), whether
\[
\exists\text{ a periodic group }G\quad\exists\text{ an injective semigroup homomorphism }S\longrightarrow\mathcal T(G)?
\]
The periodic group in this second formulation need not be finite. The supplied wording does not explicitly explain the relationship intended between these two questions; no equivalence between them is assumed here.

For the background comparison, a semigroup \(S\) is a divisor of a semigroup \(T\) if there exist a subsemigroup \(U\subseteq T\) and a surjective semigroup homomorphism \(U\to S\). The source states that not every finite semigroup has the embedding property, without separating the finite-group and periodic-group variants. It further states that every semigroup is a divisor of a power semigroup over a periodic group, citing [BersKubl], and reports a positive answer for the divisor question over finite groups, citing [Henckell]. These are supplied background claims, not independently verified assertions; divisibility and embeddability are distinct conditions.
\subsection{Short English statement}
The problem asks whether one can decide from a finite semigroup's multiplication table whether it embeds into a power semigroup of some finite group, and apparently also whether it embeds into one over some periodic group. The precise power-semigroup convention and the interpretation of the parenthetical variant require source clarification. The registry labels the problem open; this draft does not independently verify that status.
\subsection{Sources}
\begin{itemize}
\item Local UnsolvedMath registry, permanent ID \(800003\), exact source problem code \texttt{catalogue entry}, title \emph{The Lyapin's Notebook --- Problem 3}. The supplied record labels the problem \texttt{open} and gives the proposed year as 2026. These are supplied registry labels, not independently verified mathematical or historical assessments.
\item Bershadsky, Kublanovsky, Mashevitzky, \emph{The Lyapin's Notebook: Unsolved Problems in Semigroup Theory} (2026), Problem 3, \url{https://arxiv.org/abs/2604.04763}. The problem itself is attributed to S. Bershadsky and S. Kublanovsky. The supplied background records access on 2026-10-16 and extraction by \texttt{arxiv\_latex}; these are supplied provenance details, not a claim of independent access or verification. The statement refers to Chapter 5 for context.
\item [BersKubl], cited in the supplied statement for the assertion that every semigroup is a divisor of a power semigroup over a periodic group. No bibliographic entry or separate URL is supplied.
\item [Henckell], cited in the supplied statement for the positive answer concerning divisors of power semigroups over finite groups. No bibliographic entry or separate URL is supplied.
\end{itemize}
\end{document}
